\section{A uniform quantum Airy transition for the regular family}
\label{sec:quantumAiry}

This section proves a statement for an explicitly specified
meromorphic Barnes integral on the one-parameter regular slice.
It does not assert a transition theorem for arbitrary six angle
parameters, a three-manifold invariant, or an identification with
the cone-metric measure. No novelty claim is made.

\subsection{The exact normalization and the result}

Take $b>0$, $Q=b+b^{-1}$, and
\[
 \alpha=\frac{\pi-\eta}{2\pi b},\quad
 \eta_{\mathrm E}=\arccos(1/3),\quad
 h=2\pi b^2,\quad \lambda=h^{1/3}.
\]
For $\eta$ close to $\eta_{\mathrm E}$ and sufficiently small $b$,
choose an upward vertical contour $C$ satisfying
$Q<\operatorname{Re}w<4\alpha$. This lies in the usual separating
interval $3\alpha<\operatorname{Re}w<\min(4\alpha,2Q)$ and puts
all three integrand arguments in their common fundamental strip.
With the integrand of~\eqref{eq:regularquantumintegrand}, namely
\[
 I_b(w,\eta)=S_b(w-3\alpha)^4
             S_b(4\alpha-w)^3S_b(2Q-w),
\]
define
\begin{equation}
 \mathcal R_b(\eta)=
 \frac{S_b(3\alpha-Q)^2}{S_b(\alpha)^6}
 \int_C I_b(w,\eta)\,\frac{\dd w}{i}.
 \label{eq:realregularsymbol}
\end{equation}
This is the equal-label symmetric formula (2.27) of
\cite[arXiv version~3]{TeschnerVartanov}, with the measure convention specified as
$\dd w/i$. If that formula's printed $\dd w$ is read literally,
its value is $i\mathcal R_b$. Keeping this constant phase explicit
avoids imposing an unstated Barnes-measure convention.
The four identical trivalent square roots multiply to the displayed
quotient, with no remaining square-root choice. Absolute convergence
and independence of the allowed vertical line are proved below.

The functional equation $S_b(x+b^{-1})=2\sin(\pi x/b)S_b(x)$
gives the exact identity
\begin{equation}
 \frac{S_b(3\alpha-Q)^2}{S_b(\alpha)^6}
 =
 \frac{1}{4\sin^2((3\eta-\pi)/(2b^2))}
 \frac{S_b(3\alpha-b)^2}{S_b(\alpha)^6}.
 \label{eq:externalpoleclearing}
\end{equation}
Accordingly define
\begin{equation}
 \widehat{\mathcal R}_b(\eta)
 =
 4\sin^2\!\left(\frac{3\eta-\pi}{2b^2}\right)\mathcal R_b(\eta)
 =
 \frac{S_b(3\alpha-b)^2}{S_b(\alpha)^6}
 \int_C I_b(w,\eta)\,\frac{\dd w}{i}.
 \label{eq:clearedregularsymbol}
\end{equation}
The right-hand side defines the continuation across the apparent
singularities in the product on the left.

\begin{proposition}[Uniform regular-family transition]
\label{prop:regularquantumAiry}
For every fixed $Z>0$, put
\[
 \eta=\eta_{\mathrm E}
       -\frac{h^{2/3}}{6\sqrt2}\,\zeta,
       \qquad |\zeta|\le Z.
\]
As $b\downarrow0$, the integral \eqref{eq:clearedregularsymbol}
satisfies, uniformly for these real $\zeta$,
\begin{equation}
 \boxed{\quad
 \widehat{\mathcal R}_b(\eta)=
 \frac34\,\frac{h^{1/3}}{b}
 \left\{
 \operatorname{Ai}(\zeta)
 -3\sqrt2\,h^{1/3}\operatorname{Ai}'(\zeta)
 +O_Z(h^{2/3})
 \right\}.
 \quad}
 \label{eq:quantumAirytheorem}
\end{equation}
The error inside braces is additive. The physical compact
hyperbolic side has $\zeta\ge0$; negative $\zeta$ describes the
analytic continuation across the Euclidean transition.
\end{proposition}

In particular the leading prefactor is
$\frac34(2\pi)^{1/3}b^{-1/3}$. The correction can equivalently
be written by replacing the Airy argument with
$\zeta-3\sqrt2\,h^{1/3}$, at an error $O_Z(h^{2/3})$.
The proof below controls the exact quantum integrand, the whole
contour, its orientation, and the first correction.

\subsection{A uniform double-sine estimate from its real integral}

Let $\Lambda$ be the analytic primitive on
$0<\operatorname{Re}z<\pi$ specified by
\[
 \Lambda'(z)=-\log(2\sin z),\qquad \Lambda(\pi/2)=0,
\]
with the logarithm real on $(0,\pi)$. On that real interval it is
the usual Lobachevsky function.

\begin{lemma}[Strip estimate]
\label{lem:stripdoubleSine}
For $z$ in a compact subset of $0<\operatorname{Re}z<1$ and
real $\kappa$ in a bounded set,
\begin{equation}
 S_b(z/b+\kappa b)
 =
 \exp\!\left(-\frac{\Lambda(\pi z)}{\pi b^2}\right)
 (2\sin\pi z)^{\kappa-1/2}
 \bigl(1+O(b^2)\bigr).
 \label{eq:uniformstripSb}
\end{equation}
The logarithmic error and its derivatives are uniform on
smaller compact subsets. For
$\varepsilon\le\operatorname{Re}z\le1-\varepsilon$,
the real part of the logarithmic error is $O_\varepsilon(b^2)$
uniformly for all $\operatorname{Im}z$.
\end{lemma}

\begin{proof}
We use $S_b(x)=\Gamma_b(x)/\Gamma_b(Q-x)$, with
$\Gamma_b(Q/2)=1$. The convergent real-integral normalization is
\[
 \log\Gamma_b(x)=\int_0^\infty\frac{\dd u}{u}
 \left\{
 \frac{e^{-xu}-e^{-Qu/2}}{(1-e^{-bu})(1-e^{-u/b})}
 -\frac{(Q-2x)^2}{8}e^{-u}-\frac{Q-2x}{2u}
 \right\},\qquad \operatorname{Re}x>0.
\]
The first term in braces is
$(Q/2-x)/u+(Q/2-x)^2/2+O(u)$ at zero, which fixes both
subtractions. In~\cite[arXiv version~3, (A.1)]{TeschnerVartanov}
the last denominator is printed as $u$; the factor $2$ above
is necessary for convergence. Subtracting the integrals at $x$
and $Q-x$ therefore gives, for $0<\operatorname{Re}x<Q$,
\[
 \log S_b(x)=\int_0^\infty\frac{\dd u}{u}
 \left\{
 \frac{\sinh((Q/2-x)u)}
      {2\sinh(bu/2)\sinh(u/(2b))}
 -\frac{Q-2x}{u}
 \right\}.
\]
This fixes a holomorphic logarithm, real on $(0,Q)$ and zero at
$Q/2$. At the centered argument $x=z/b+b/2$, substitute
$u=2bs$ and separate the term $1/(b^2s)$ from
$\operatorname{csch}(b^2s)$. The result is
\begin{equation}
 \log S_b(z/b+b/2)
 =-\frac{\Lambda(\pi z)}{\pi b^2}+D_b(z),
 \label{eq:stripSbexactremainder}
\end{equation}
where
\begin{equation}
 D_b(z)=\frac12\int_0^\infty
 \frac{\sinh((1-2z)s)}{s\sinh s}
 \left(\operatorname{csch}(b^2s)-\frac1{b^2s}\right)\dd s.
 \label{eq:Sbremainderintegral}
\end{equation}
Here the leading integral is the elementary identity
\[
 \frac12\int_0^\infty
 \left\{\frac{\sinh((1-2z)s)}{\sinh s}-(1-2z)\right\}
 \frac{\dd s}{s^2}=-\frac{\Lambda(\pi z)}{\pi}.
\]
For example, differentiating twice gives $\pi\cot(\pi z)$
by the partial fraction expansion of the cotangent. Its first
derivative at $1/2$ is
$\int_0^\infty(s^{-2}-(s\sinh s)^{-1})\dd s=\log2$,
as follows from $\operatorname{csch}s=\coth(s/2)-\coth s$
and a change of scale. Its value at $1/2$ is zero. These three
facts determine the primitive and its constant.
For $u>0$,
$|\operatorname{csch}u-u^{-1}|\le C u$.
Near $s=0$ the numerator $\sinh((1-2z)s)$ is $O(s)$,
and at infinity the strip condition supplies exponential decay.
This proves $D_b=O(b^2)$ on compact subsets. Holomorphy and
Cauchy's estimate give the derivative statement.

For the global real-part bound use
\[
 \operatorname{Re}\sinh((1-2z)s)
 =\sinh((1-2\operatorname{Re}z)s)
       \cos(2\operatorname{Im}z\,s).
 \]
The same integrable bound is then independent of
$\operatorname{Im}z$. Finally replace $z$ in
\eqref{eq:stripSbexactremainder} by
$z+(\kappa-\frac12)b^2$ and Taylor expand the leading
primitive. Its second and higher derivatives are bounded on
closed substrips, including as $|\operatorname{Im}z|\to\infty$.
This gives \eqref{eq:uniformstripSb} and the corresponding
global modulus estimate.
The centered variable is $z=bx-b^2/2$. Its strip
$0<\operatorname{Re}z<1$ is narrower than the fundamental
$x$-strip; no uniform small-$b$ estimate is asserted up to the
boundary of that larger strip.
\end{proof}

\subsection{The smooth quantum phase and amplitude}

Use the corrected variable $v=2\pi bw$, and set
\begin{equation}
 \begin{aligned}
 \theta_1&=\tfrac12(v-3\pi+3\eta),&
 \theta_2&=\tfrac12(4\pi-4\eta-v),\\
 \theta_3&=\tfrac12(4\pi-v),&
 s&=\tfrac12(\pi-\eta).
 \end{aligned}
 \label{eq:regularstripangles}
\end{equation}
The three integrand arguments are
$\theta_1/(\pi b)$, $\theta_2/(\pi b)$, and
$\theta_3/(\pi b)+2b$.
The cleared prefactor uses $3s/(\pi b)-b$ and $s/(\pi b)$.
All five scaled arguments are in a closed substrip near the fold.
Lemma~\ref{lem:stripdoubleSine} therefore gives
\begin{equation}
 \frac{S_b(3\alpha-b)^2}{S_b(\alpha)^6}I_b(w,\eta)
 =
 \mathcal A(v,\eta)\,e^{G(v,\eta)/h}
 \bigl(1+O(h)\bigr)
 \label{eq:fullsmoothquantumphase}
\end{equation}
on compact subsets, where
\begin{align}
 G(v,\eta)
 &=-2\{4\Lambda(\theta_1)+3\Lambda(\theta_2)
                              +\Lambda(\theta_3)\}
        +12\Lambda(s)-4\Lambda(3s),\notag\\
 \mathcal A(v,\eta)
 &=(2\sin\theta_1)^{-2}(2\sin\theta_2)^{-3/2}
   (2\sin\theta_3)^{3/2}
   \left(\frac{\sin s}{\sin3s}\right)^3.
 \label{eq:clearedphaseamplitude}
\end{align}
The powers are continued from their positive real values.
For modulus estimates on the full vertical line, the same formula
holds with a bounded multiplicative error, uniformly in its
imaginary coordinate; this follows from the real-part estimate
in the lemma.

At $v_*=4\pi-5\eta_{\mathrm E}$ the exact values are
\[
 \begin{array}{c|ccc}
 &\theta_1&\theta_2&\theta_3\\ \hline
 \sin\theta_j&1/3&1/\sqrt3&1/(9\sqrt3)\\
 \cot\theta_j&2\sqrt2&\sqrt2&-11\sqrt2
 \end{array}
 \qquad
 \frac{\sin s}{\sin3s}=3.
\]
Consequently
\begin{equation}
 \mathcal A_*=\frac94,\qquad
 (\log\mathcal A)_v\big|_*=7\sqrt2.
 \label{eq:regularAiryamplitude}
\end{equation}
The phase has the following fold data; its value is justified below,
and its derivatives follow from~\eqref{eq:clearedphaseamplitude}:
\begin{equation}
 \begin{gathered}
 G_*=G_v|_*=G_{vv}|_*=G_\eta|_*=0,\\
 G_{vvv}|_*=54,\qquad G_{v\eta}|_*=18\sqrt2,\qquad
 G_{\eta\eta}|_*=36\sqrt2,\\
 G_{vvvv}|_*=-648\sqrt2,\qquad G_{vv\eta}|_*=-18.
 \end{gathered}
 \label{eq:regularAiryderivatives}
\end{equation}
The value $G_*=0$ requires fixing the constant, not just
differentiating the phase. Use the real Lobachevsky form of
\cite[Section~1.2, (1.14)--(1.18)]{MurakamiYano}, extending
the real function by oddness and period $\pi$. For equal angles
its potential and correction are
\[
 \begin{split}
 U_\Lambda(v,\eta)
 &=\Lambda(v/2)+3\Lambda(v/2+2\eta-\pi)
       -4\Lambda((v+3\eta-\pi)/2),\\
 \Delta_\Lambda(\eta)
 &=2\Lambda((3\eta-\pi)/2)-6\Lambda((\eta-\pi)/2).
 \end{split}
\]
Periodicity and oddness give, term by term,
\[
 \begin{split}
 U_\Lambda&=-4\Lambda(\theta_1)-3\Lambda(\theta_2)
                         -\Lambda(\theta_3),\\
 \Delta_\Lambda&=6\Lambda(s)-2\Lambda(3s),\qquad
 G=2(U_\Lambda+\Delta_\Lambda).
 \end{split}
\]
At the two unit-circle saddle roots on the compact hyperbolic
side, Theorem~2 of that paper, in this real formulation, says
$|U_\Lambda+\Delta_\Lambda|=W$. Both roots converge to the
lift $v_*$, while the regular tetrahedral volume tends to zero
by~\eqref{eq:regularvolumebyarea}. Continuity proves $G_*=0$.
This argument fixes the real constant without selecting
unspecified branches of a complex dilogarithm.

\subsection{The full contour and its tail estimate}

The poles of the first double-sine factor are at
$w=3\alpha-mb-n/b$; the other two factors have their poles
at $w=4\alpha+mb+n/b$ and $w=2Q+mb+n/b$.
Thus the separating interval remains nonempty for
$0<\alpha<2Q/3$. Decreasing the equal real labels from
$\alpha=Q/2$ to the compact values preserves this interval.
Along this path $\alpha>Q/4$, so one may choose
$\max(3\alpha,Q)<\operatorname{Re}w<4\alpha\le2Q$ and keep
all three arguments in the common strip throughout.
The integral has no pinch on this path.

First justify moving an arbitrary defining contour to the fold
neighborhood at each fixed $b$. The vertical asymptotics
\cite[(A.11) and (A.13)]{TeschnerVartanov} give
\[
 |S_b(a+iT)|=e^{\pi(a-Q/2)|T|}(1+o(1)),\qquad |T|\to\infty,
\]
uniformly for $a$ in a fixed compact real interval. The sum of
the coefficients in the eight factors of $I_b$ is
\[
 4(w_R-3\alpha-Q/2)+3(4\alpha-w_R-Q/2)
                       +(2Q-w_R-Q/2)=-2Q.
\]
Consequently $|I_b(w_R+iT,\eta)|=O(e^{-2\pi Q|T|})$,
uniformly between any two fixed separating lines. Horizontal
closing segments vanish and no poles are crossed. This proves
absolute convergence, contour independence, and agreement with
the continued contour in (2.27). The same bound is locally
uniform for complex $\alpha$ near the real continuation path,
so the integral is holomorphic there. The external factor is
meromorphic; \eqref{eq:externalpoleclearing} removes its poles
on this path. In a neighborhood of the fold, both remaining
prefactor arguments are in $(0,Q)$, and the cleared expression
is holomorphic and real for real $\eta$.

Only after this fixed-$b$ relocation do we apply the uniform
centered-strip estimates near $v_*$. They apply to the entire
vertical line there, including its unbounded imaginary part.

On the vertical line $v=v_*+iy$ at $\eta=\eta_{\mathrm E}$,
put $r=\tanh(y/2)$. For $y>0$, direct differentiation gives
\begin{equation}
 -\frac{\dd}{\dd y}\operatorname{Re}G(v_*+iy,\eta_{\mathrm E})
 =\Psi(r),
 \label{eq:verticalphasegap}
\end{equation}
where
\[
 \Psi(r)=4\arctan(2\sqrt2r)
          +3\arctan(\sqrt2r)-\arctan(11\sqrt2r).
\]
Here $\Psi(0)=0$ and
\begin{equation}
 \Psi'(r)=
 \frac{864\sqrt2\,r^2(3+11r^2)}
 {(1+8r^2)(1+2r^2)(1+242r^2)}>0
 \qquad(r>0).
 \label{eq:verticalphasepositivity}
\end{equation}
It follows that the fold is the unique maximum of the real part
of $G$ on this vertical line. In addition,
\[
 \operatorname{Re}G(v_*+iy,\eta_{\mathrm E})
   =-27\sqrt2\,y^4+O(y^6),\qquad
 \Psi(\tanh(y/2))\longrightarrow2\pi
 \quad(y\to+\infty).
\]
Complex conjugation gives the corresponding estimates for
negative $y$.

These bounds persist in the form needed below under small
changes of $\eta$ and of the real contour coordinate. On bounded
annuli in $y$ this follows by continuity. For large positive $y$,
write $a_j=\theta_j(v_R,\eta)\in(0,\pi)$. The arguments of
$2\sin(a_1+iy/2)$ and $2\sin(a_j-iy/2)$, $j=2,3$,
are respectively $\pi/2-a_1+O(e^{-y})$ and
$a_j-\pi/2+O(e^{-y})$. Since
$4a_1+3a_2+a_3=2\pi$, it follows that
$\partial_y\operatorname{Re}G=-2\pi+O(e^{-y})$,
uniformly for nearby real $v_R,\eta$. Conjugation treats the
negative tail. The amplitude in
\eqref{eq:clearedphaseamplitude} is $O(e^{-|y|})$ there.
The global modulus bound from
Lemma~\ref{lem:stripdoubleSine} transfers these estimates to
the exact integrand, so all tails outside a fixed fold
neighborhood are exponentially small in $h^{-1}$. More
explicitly, for each sufficiently small fixed $\delta>0$ there
are $m_\delta>0$ and $Y_0>\delta$, uniform in nearby real
parameters, such that
\[
 \operatorname{Re}G(v_R+iy,\eta)
 \le -m_\delta-\pi(|y|-Y_0)_+,
 \qquad |y|\ge\delta.
\]

\subsection{Airy localization and the first correction}

\begin{proof}[Proof of Proposition~\ref{prop:regularquantumAiry}]
Write $t=\eta_{\mathrm E}-\eta$ and $\lambda=h^{1/3}$.
Choose a fixed $c>0$ and move the contour to
\[
 v=v_*+c\lambda+iy.
\]
For all sufficiently small $b$ this deformation stays within
the common pole-free strip. The tail bounds just proved make
the horizontal closing segments vanish, so no residues or
undetermined cycle coefficients enter.

At this displaced real point,
$G_{vv}=54c\lambda+O_Z(\lambda^2)>0$ and
$G_{vvvv}=-648\sqrt2+O_Z(\lambda)<0$.
Taylor expansion in $iy$, with real derivatives at the real
point, has real part
$-G_{vv}y^2/2+G_{vvvv}y^4/24+O(y^6)$. It gives constants
$c_1,c_2,\delta>0$ such that
\begin{equation}
 \begin{split}
 &\operatorname{Re}\{G(v_*+c\lambda+iy,\eta)
                     -G(v_*+c\lambda,\eta)\}\\
 &\hspace{25mm}\le -c_1\lambda y^2-c_2y^4,
       \qquad |y|\le\delta.
 \end{split}
 \label{eq:AiryGaussianbound}
\end{equation}
The constants can be taken uniform for $|\zeta|\le Z$.
Outside this neighborhood use the previously established
fixed tail gap and linear decay.

Now set
\[
 x=v-v_*=\frac{\lambda}{3}Y,\qquad
 t=\frac{\lambda^2}{6\sqrt2}\zeta.
\]
The rescaled contour is upward with $\operatorname{Re}Y=3c$.
The derivative table gives, through weighted order four,
\begin{align}
 G(v_*+x,\eta_{\mathrm E}-t)
 &=9x^3-18\sqrt2\,tx
   +18\sqrt2\,t^2-27\sqrt2\,x^4+9tx^2
   +R_G(x,t),\notag\\
 \frac{\mathcal A(v_*+x,\eta_{\mathrm E}-t)}{9/4}
 &=1+7\sqrt2\,x+O(x^2+|t|).
 \label{eq:Airyweightedexpansion}
\end{align}
Analyticity gives the explicit local bound
\[
 |R_G(x,t)|\le C\bigl(|x|^5+|t||x|^3+t^2|x|+|t|^3\bigr).
\]
The exact quantum error in
\eqref{eq:fullsmoothquantumphase} is $O(h)=O(\lambda^3)$
and hence does not affect the first correction.
Consequently, the rescaled integrand is
\begin{equation}
 \frac94\,e^{Y^3/3-\zeta Y}
 \left\{1+\lambda P(Y,\zeta)+O_Z(\lambda^2)\right\},
 \label{eq:Airyscaledquantum}
\end{equation}
where the pointwise remainder is locally uniform and
\begin{equation}
 P(Y,\zeta)=
 \frac{7\sqrt2}{3}Y-\frac{\sqrt2}{3}Y^4
 +\frac{\zeta}{6\sqrt2}Y^2
 +\frac{\sqrt2}{4}\zeta^2.
 \label{eq:Airycorrectionpolynomial}
\end{equation}

For completeness, the integrated remainder is also
$O_Z(\lambda^2)$. On
$|Y|\le\lambda^{-1/12}$, analytic Taylor remainders in
\eqref{eq:Airyweightedexpansion}, followed by exponentiation,
give an error bounded by
$C_Z\lambda^2(1+|Y|^8)$, against the Gaussian
$|e^{Y^3/3-\zeta Y}|\le C_Ze^{-3c(\operatorname{Im}Y)^2}$.
Indeed $|R_G|/h\le C_Z\lambda^2(1+|Y|^5)$, and the
square of the order-$\lambda$ phase correction has degree
eight. The growing cutoff makes both exponent remainders
uniformly small.
Equation~\eqref{eq:AiryGaussianbound} supplies a uniform
Gaussian majorant for the exact rescaled integrand, since
$G(v_*+c\lambda,\eta)/h=O_Z(1)$. The part with
$\lambda^{-1/12}<|Y|$ inside the fixed fold neighborhood
is smaller than every power of $\lambda$; the global contour
tail is exponentially small in $h^{-1}$. This justifies
integration of the expansion and its first correction.

With the chosen orientation,
\[
 \frac1{2\pi i}\int_{\operatorname{Re}Y=3c}
       e^{Y^3/3-\zeta Y}\,\dd Y=\operatorname{Ai}(\zeta);
\]
this follows from~\cite[(9.5.4)]{DLMF} with its contour
oriented from the lower ray to the upper ray. To deform those
rays to $\operatorname{Re}Y=3c$, close at imaginary heights
$\pm T$. On each horizontal connector, with
$3c\le x=\operatorname{Re}Y\le T/\sqrt3$,
\[
 \operatorname{Re}(Y^3/3-\zeta Y)
 \le-\tfrac89xT^2+Zx,
\]
so the connector integrals vanish. Inserting $Y^k$ yields
$(-1)^k\operatorname{Ai}^{(k)}(\zeta)$.
Using
\[
 \operatorname{Ai}''(\zeta)=\zeta\operatorname{Ai}(\zeta),
 \qquad
 \operatorname{Ai}^{(4)}(\zeta)
       =2\operatorname{Ai}'(\zeta)
          +\zeta^2\operatorname{Ai}(\zeta)
\]
in \eqref{eq:Airycorrectionpolynomial}, all terms proportional
to $\zeta^2\operatorname{Ai}$ cancel. The result is
$-3\sqrt2\,\operatorname{Ai}'(\zeta)$.
Finally,
$\dd w/i=\lambda\,\dd Y/(6\pi i b)$.
Multiplication by $9/4$ gives the coefficient
$3\lambda/(4b)$ in \eqref{eq:quantumAirytheorem}.
\end{proof}

\subsection{Numerical evidence and reproducibility status}

The values below come from a new evaluation of the exact cleared
integral~\eqref{eq:clearedregularsymbol}, implemented in
\path{experiments/section15_quantum/}. Write
\[
 Q_b(\zeta)=\frac{4b}{3h^{1/3}}\widehat{\mathcal R}_b(\eta),
 \qquad
 Q_b^{(1)}(\zeta)=\operatorname{Ai}(\zeta)
       -3\sqrt2\,h^{1/3}\operatorname{Ai}'(\zeta).
\]
The evaluator uses the exact centered identity
\eqref{eq:stripSbexactremainder}. A finite expansion of $D_b$
in derivatives of the cotangent has an explicit rational
remainder bound uniform in a declared centered strip. This
approximates the double-sine correction, not the Airy phase.
All five logarithmic terms, with weights $4,3,1,2,-6$, are
combined before exponentiation. The logarithms and the
Lobachevsky primitive have the branches fixed above.

The outer Gauss--Legendre quadrature uses the upward contour
\[
 v=v_*+\frac{\lambda}{3}(C+it),\qquad C>0,
\]
with measure $\dd w/i$; thus $C=3c$ in the proof.
Each of the seven parameter pairs described below is evaluated
at 60 decimal digits with $C=1$, $|t|\le12$, then at 90 digits
with $C=1$, $|t|\le14$, and again at 90 digits with $C=5/4$
and two additional correction terms. The contribution beyond
a shorter cutoff is recorded separately. The table rounds the
90-digit, $C=1$ values to six decimals.
\begin{center}
\begin{tabular}{rrrr}
\toprule
$b$&$\zeta$&Computed $Q_b$&Airy $Q_b^{(1)}$\\
\midrule
$0.0125$&$0$&$0.459086$&$0.464164$\\
$0.00625$&$0$&$0.422572$&$0.423779$\\
$0.003125$&$0$&$0.398047$&$0.398339$\\
$0.0125$&$1$&$0.214333$&$0.202400$\\
$0.00625$&$1$&$0.182316$&$0.177567$\\
$0.003125$&$1$&$0.163805$&$0.161924$\\
\bottomrule
\end{tabular}
\end{center}
For the three $\zeta=1$ rows, the newly computed values of
$(Q_b-Q_b^{(1)})/h^{2/3}$ are approximately
$1.208081684$, $1.211366819$, and $1.208878763$.
A negative transition parameter is also checked:
$Q_{1/160}(-1)\approx0.517966671856531$.
The largest observed 60/90-digit difference in these seven
values is less than $4\times10^{-57}$; changing the contour
and correction order changes them by less than
$1.3\times10^{-91}$. These are comparisons between numerical
runs, not guaranteed numbers of correct digits.

Direct quadrature of $D_b$ and of the convergent real
$\log S_b$ integral, together with the reflection and shift
identities, checks the special-function evaluation independently
of the correction expansion. A separate computation integrates
the combined Barnes logarithm directly, using QUADPACK in
ordinary double precision; its seven normalized values agree
with the multiprecision calculation within $2\times10^{-12}$.
The rational bound controls correction-series truncation only.
Its propagation uses a numerical absolute-integral estimate;
the outer quadrature errors, sign tests, and tail-majorant
formulas are evaluated without outward rounding. No total
interval enclosure of the quantum integral is claimed.

The supplied draft's evaluator and raw output remain
unrecovered. Its six reported values are preserved in
\path{experiments/section15/historical-table.json}.
The new first value is
$0.459086488736078952\ldots$, replacing the previously printed
$0.459087$ by $0.459086$; the original rounding procedure is
unknown. The other five values agree after rounding. This is
new computational evidence with its own provenance, not a
recovery of the missing original calculation.

The earlier companion \path{experiments/section15/check.py}
checks the finite fold derivatives, the rational identity for
$\Psi'$, and the first-correction cancellation in exact
arithmetic with $\sqrt2$ adjoined. It also checks the Airy
column at 80 and 120 digits against a separately coded ODE
series~\cite[Sections~9.2 and~9.4]{DLMF}; its existing result
remains a frozen record of those finite checks.
Reproduction instructions for the new evaluator are in
\path{experiments/section15_quantum/README.md}; its full
result is written to \path{artifacts/section15_quantum/result.json}.
Source identities, independent reviews, and generation records
are collected in
\path{docs/execution/section15-analytic-review-2026-11-01/README.md}.
The analytic proof of Proposition~\ref{prop:regularquantumAiry}
and this finite numerical evidence remain distinct.

\subsection{Interpretation and scope}

The correction to the earlier pole argument is substantive.
The exact quantum integrand is smooth at the correct fold;
the remaining real-$b$ resonances are an explicit scalar
factor outside the integral. For example, at $\eta=\eta_{\mathrm E}$
the scalar factor has candidate resonances along
\[
 b^2=\frac{3\eta_{\mathrm E}-\pi}{2\pi m},
       \qquad m=1,2,\ldots.
\]
The numerator in \eqref{eq:clearedregularsymbol} is nonzero
for sufficiently small $b$, uniformly in each fixed physical
window $0\le\zeta\le Z$: the continuous positive function
$\operatorname{Ai}(\zeta)$ has a positive minimum on $[0,Z]$,
and the additive correction and error tend uniformly to zero.
Thus zeros of the squared sine give genuine double poles of
the original normalization in this window, including all
sufficiently large $m$ in the displayed sequence. For negative
$\zeta$, the additive asymptotic remains uniform even at
zeros of $\operatorname{Ai}$; no uniform noncancellation or
relative-error assertion is made there.
The cleared quantity has the uniform asymptotic in the
proposition even through these resonances.

In geometric terms,
\[
 \zeta=\frac{A}{2\sqrt3\,h^{2/3}}
          +O\!\left(\frac{A^2}{h^{2/3}}\right),
 \qquad A\asymp b^{4/3}.
\]
The decaying Airy branch selected by the contour agrees with
the decaying volume exponential: its classical saddle
exponent is $-2W/h=-W/(\pi b^2)$.
The proposition is uniform for bounded $\zeta$; it does not,
by itself, justify taking $\zeta\to\infty$ simultaneously
with $b\to0$ or keeping a nonzero tetrahedral volume fixed.
Those matching limits require additional estimates.

This proves a quantum transition for one specified regular
family, beyond the classical cubic calculation. Extension
to arbitrary nondegenerate Euclidean tetrahedral shapes,
comparison of normalizations used in a global state integral,
and gluing into a topological invariant remain separate
problems. In particular no new case of the colored-Jones
volume conjecture, and no identification of the integrated
quantum measure with the cone-metric measure, follows from
this proposition alone.

